

\chapter{Függelékek és production checklist}

\section{Aktuális verifikációs megjegyzés}

A jelenlegi Velran forrásfa teljesítette a repository-szintű Ellenőrzött fejlesztési mérföldkövet: a formázás, a teljes workspace tesztkészlet, a \code{./verify.sh} és a helyi tesztkiszolgálás sikeresen lefutott. Production release előtt az egzakt release-forrásfán ismét futtasd a lock/provenance frissítést és a repository teljes verifikációját, majd gyűjtsd össze a célkörnyezet deployment-, recovery- és operátori bizonyítékait. A repository verifikáció szükséges, de nem helyettesíti a production-környezeti release evidence-et.

\noindent\textbf{Tanulási út: Gyors út: végső release gate.} A checklisttel bizonyítsd a release readiness állapotot, ne emlékezetből feltételezd.\par\medskip

Ez a fejezet nem új nyelvi elemeket vezet be. Az előző fejezetekből összeálló, röviden használható üzemeltetési segédlet: tipikus könyvtárstruktúra, reverse proxy minták, systemd unit, logrotate és egy release előtti ellenőrzőlista. A minták kiindulópontok; a konkrét domaineket, tanúsítványokat, jogosultságokat és erőforráslimiteket a saját környezethez kell igazítani.

\section{Ajánlott production könyvtárstruktúra}

A Velran helyben telepített szoftver, ezért a könyv következetesen a \code{/usr/local} hierarchiát használja a programhoz és a konfigurációhoz. Az alkalmazás állapota és a logok ettől külön élnek.

\begin{lstlisting}
/usr/local/bin/
  velran-server
  velran-cli

/usr/local/etc/velran/
  server.toml
  rate-limits.toml
  resource-profiles.toml

/run/secrets/velran/
  database-url
  redis-url
  local-auth-database-url

/srv/velran/
  current/                 # current application release
  releases/
    2026-09-04.1/
  data/                    # AppFs / durable application data

/var/cache/velran/
  native/                  # verified native cdylib cache

/var/log/velran/
  server.log
  access.log
  audit.log
\end{lstlisting}

Az immutable release tartalom legyen lehetőleg read-only. Tranzakciós rolling-source oldalnál külön alkalmazás-forráskönyvtárat használj, amelyet a deployment/feltöltési workflow birtokol és a Velran service olvasni tud; maga az alkalmazás-runtime ne kapjon általános írási jogot a saját forrásfájaira. Mindkét modellben csak ott legyen szerveroldali írási jog, ahol tartós állapot, log vagy explicit kezelt natív cache ezt igényli. A secret fájl ne legyen az alkalmazás forrás/release könyvtárában és ne kerüljön verziókezelésbe.

\section{Apache reverse proxy teljes minimálminta}

A példa egyetlen Apache reverse proxyt feltételez. A publikus TLS az Apache feladata, a Velran pedig csak loopbacken figyel.

Az ehhez tartozó Velran minimum:

\begin{lstlisting}
[server]
listen = "127.0.0.1:8080"
insecure_dev_cookies = false

[tls]
public_host = "www.example.com"

[web]
trusted_proxy_cidrs = ["127.0.0.1/32", "::1/128"]
\end{lstlisting}

Reverse proxy production módban a certificate/key a proxyn marad, de a trusted \code{public_host} kell a Host/Origin pinninghez. A proxy feladata a hiteles \code{X-Forwarded-Proto: https} metadata előállítása.

\begin{lstlisting}
<VirtualHost *:80>
    ServerName www.example.com
    Redirect permanent / https://www.example.com/
</VirtualHost>

<VirtualHost *:443>
    ServerName www.example.com

    SSLEngine on
    SSLCertificateFile \
      /etc/letsencrypt/live/www.example.com/fullchain.pem
    SSLCertificateKeyFile \
      /etc/letsencrypt/live/www.example.com/privkey.pem

    ProxyPreserveHost On
    ProxyPass        / http://127.0.0.1:8080/
    ProxyPassReverse / http://127.0.0.1:8080/

    RequestHeader set X-Forwarded-Proto "https"
</VirtualHost>
\end{lstlisting}

Tipikus Apache modulok: \code{proxy}, \code{proxy_http}, \code{headers} és \code{ssl}. A Velran trusted-proxy listája csak a tényleges proxyt foglalja magában. Loopback proxy esetén ez tipikusan \code{127.0.0.1/32} és szükség esetén \code{::1/128}. Internetméretű trusted proxy tartomány hibás konfiguráció.

\section{Nginx reverse proxy teljes minimálminta}

\begin{lstlisting}
server {
    listen 80;
    server_name www.example.com;
    return 301 https://$host$request_uri;
}

server {
    listen 443 ssl;
    server_name www.example.com;

    ssl_certificate
      /etc/letsencrypt/live/www.example.com/fullchain.pem;
    ssl_certificate_key
      /etc/letsencrypt/live/www.example.com/privkey.pem;

    location / {
        proxy_pass http://127.0.0.1:8080;
        proxy_http_version 1.1;
        proxy_set_header Host $host;
        proxy_set_header X-Forwarded-Proto https;
        proxy_set_header X-Forwarded-For $remote_addr;
    }
}
\end{lstlisting}

A példában az Nginx közvetlenül látja a klienst, ezért a forwarding címet maga állítja elő. Ne örökíts tovább internet felől érkező tetszőleges \code{X-Forwarded-*} értéket authorityként. Több proxy esetén a trust láncot explicit kell megtervezni.

\section{systemd service minta}

A repository konzervatív baseline mintája: \code{examples/systemd/velran.service}. Az alábbi hardened forma ezen felül megmutatja a perzisztens natív cache-hez szükséges írható útvonalat is:

\begin{lstlisting}
[Unit]
Description=Velran application
Wants=network-online.target
After=network-online.target

[Service]
Type=simple
User=velran
Group=velran
WorkingDirectory=/srv/velran/current
ExecStartPre=/usr/local/bin/velran-server \
  --config /usr/local/etc/velran/server.toml \
  --check-config
ExecStart=/usr/local/bin/velran-server \
  --config /usr/local/etc/velran/server.toml
ExecReload=/bin/kill -HUP $MAINPID
Restart=on-failure
RestartSec=2s
TimeoutStopSec=45s

MemoryMax=1G
MemorySwapMax=0
CPUQuota=200%
TasksMax=256
LimitNOFILE=8192
# Direct TLS 80/443 eseten add hozza:
# AmbientCapabilities=CAP_NET_BIND_SERVICE
# CapabilityBoundingSet=CAP_NET_BIND_SERVICE
NoNewPrivileges=yes
PrivateTmp=yes
PrivateDevices=yes
ProtectSystem=strict
ProtectHome=yes
ProtectKernelTunables=yes
ProtectKernelModules=yes
ProtectControlGroups=yes
RestrictSUIDSGID=yes
LockPersonality=yes
ReadWritePaths=/srv/velran/data /var/cache/velran/native /var/log/velran

[Install]
WantedBy=multi-user.target
\end{lstlisting}

A minta a perzisztens natív cache könyvtárat is írhatóként engedi. Ha a \code{native.cache_dir} máshová mutat, a systemd filesystem sandboxot ugyanahhoz az explicit könyvtárhoz igazítsd; ne oldd fel általánosan a \code{ProtectSystem=strict} védelmet. Rolling-source telepítésnél a feltöltési workflow írja az alkalmazás-forráskönyvtárat, a Velran service számára olvasási jog elegendő.\par\medskip

A systemd hard limit és a Velran request/runtime resource policyja egymást kiegészítő védelmi réteg. A systemd baseline maga legyen a process cgroup authority; ilyenkor ne kapcsold be mellette a Velran \code{[cgroup]} író módját. Közvetlen 80/443 bindnál csak a \code{CAP_NET_BIND_SERVICE} capabilityt add, reverse proxy mögötti magas portnál ezt is hagyd el. A \code{TimeoutStopSec} legyen nagyobb, mint az alkalmazott graceful shutdown időablak. Behind-proxy módban a SIGHUP tranzakciósan újratöltheti a domain/application hosting állapotot, míg process-szintű konfigurációváltozás után továbbra is kontrollált restart kell; pusztán alkalmazásforrás-változást normál esetben az automatikus source-reload supervisor aktivál.

\section{logrotate minta}

\begin{lstlisting}
/var/log/velran/server.log /var/log/velran/access.log {
    daily
    rotate 14
    compress
    delaycompress
    missingok
    notifempty
    create 0640 velran velran
    sharedscripts
    postrotate
        systemctl kill -s HUP --kill-who=main \
          velran.service >/dev/null 2>&1 || true
    endscript
}

/var/log/velran/audit.log {
    daily
    rotate 90
    compress
    delaycompress
    missingok
    notifempty
    create 0640 velran velran
    sharedscripts
    postrotate
        systemctl kill -s HUP --kill-who=main \
          velran.service >/dev/null 2>&1 || true
    endscript
}
\end{lstlisting}

A konkrét retention jogi és kockázati döntés. Az auditlog általában hosszabb megőrzést és központi továbbítást érdemel; a minta számai nem univerzális előírások.

\section{Release előtti production gate}

A következő lista szándékosan rövid és végrehajtható. A részletes indoklást az előző fejezetek adják.

\begin{enumerate}
  \item \textbf{Forrás és build:} a release commit azonosított; \code{Cargo.lock} rögzített; a platform buildje \code{cargo build --locked --release -p velran-server -p velran-cli}; az alkalmazás forrása \code{velran-cli check} ellenőrzésen átment.
  \item \textbf{Migration:} production előtt \code{migrate status} és \code{migrate verify}; a jóváhagyott schema változás explicit \code{migrate apply}; utána újabb verify.
  \item \textbf{Konfiguráció:} a valódi production \code{server.toml} config-checken átment. Az effective config redaktált kimenete reviewzható.
  \item \textbf{Secrets:} DB/Redis/auth secret fájlok nem a repositoryban vannak, szűk filesystem jogosultsággal rendelkeznek, és nem kerülnek logba vagy release-recordba.
  \item \textbf{HTTP edge:} a domain, TLS és Host kezelés helyes. A trusted-proxy lista kizárólag a tényleges proxykat tartalmazza; a certificate-renewal felelőse és a lejárat monitorozása explicit.
  \item \textbf{Auth:} production cookie policy aktív; local auth/LDAP konfiguráció ellenőrzött; admin/MFA recovery folyamat dokumentált.
  \item \textbf{CSRF és CORS:} state-changing böngészős kérés CSRF-védett; CORS allowlist explicit; credentialed CORS esetén nincs wildcard origin.
  \item \textbf{Adatbázis:} az alkalmazás DB user least privilege; pool/budget összhangban van a backend kapacitásával; fontos queryk indexelése ismert.
  \item \textbf{Storage/upload:} AppFs root confinement és jogosultságok helyesek; upload byte/pixel limitek aktívak; release tartalom nem írható alkalmazásból.
  \item \textbf{Redis/cache/rate limit:} multi-instance környezetben a közös security/runtime state Redisben van; Redis kiesési viselkedése ismert és tesztelt.
  \item \textbf{Egress:} csak named outbound target engedélyezett; host/CIDR/port/TLS/timeout/méret policy explicit; secret nem kerül URL-be vagy logba.
  \item \textbf{Erőforrások/host:} request/runtime limitek és process hard limitek konzisztensen be vannak állítva; cgroup authorityként systemd vagy külön delegált Velran cgroup legyen, ne mindkettő; legyen elegendő lemez/inode tartalék és működő host időszinkron.
  \item \textbf{Observability:} server/access/audit log írható; request ID működik; metrics és readiness elérhető a megfelelő belső hálózatról; riasztási minimum definiált.
  \item \textbf{Backup/recovery:} van friss, ellenőrzött DB + AppFs + identity backup; recovery manifest rendelkezésre áll; restore drill nem csak papíron létezik.
  \item \textbf{Deploy/rollback:} immutable deploymentnél ismert a release artifact/checksum és a rollback cél; rolling-source deploymentnél ismert az aktív source/generation identity, és bizonyított, hogy az elutasított candidate mellett az előző generation marad aktív. Schema-változásnál tisztázott, hogy app rollback vagy teljes restore lehetséges.
  \item \textbf{Indítás utáni evidence:} liveness, readiness és legalább egy üzleti smoke kérés sikeres; logban nincs startup fallback vagy váratlan policy-hiba.
\end{enumerate}

\subsection*{Platform-release security evidence}

A Velran platform release csak akkor nevezhető release-ready állapotúnak, ha pontosan az a source tree átmegy ezen:

\begin{lstlisting}
cargo fmt --all -- --check
cargo metadata --locked --format-version 1 >/dev/null
cargo check --workspace --locked
cargo test --workspace --locked
./verify.sh
\end{lstlisting}

A supply-chain envelope egyezzen a \code{Cargo.lock} és \path{SUPPLY-CHAIN-CAPABILITIES.txt} állapottal. Release packaging nem generálhat vagy pecsételhet le Cargo szerint stale lockfile-t. Minden dokumentáció/book változás után a release manifestet is újra kell verifikálni.

\section{Gyors hibadiagnosztikai térkép}

\begin{longtable}{L{0.23\textwidth}L{0.30\textwidth}L{0.38\textwidth}}
\toprule
Tünet & Első hely, ahol nézd & Következő kérdés \\
\midrule
\endhead
Nem indul a service & config-check, server log & Config/secret/path/policy vagy dependency readiness? \\
\addlinespace
502/504 a proxyban & proxy log, Velran listener/readiness & Fut-e a processz, loopback port helyes-e, readiness zöld-e? \\
\addlinespace
400/415 & access log + route body-contract & Malformed input, rossz schema vagy rossz Content-Type érkezett? \\
\addlinespace
401 & auth/session audit esemény & Nincs session, lejárt, invalidálva vagy MFA hiányzik? \\
\addlinespace
403 & route/object authorization, CSRF/CORS & Ki az actor, melyik policy tagadta meg? \\
\addlinespace
404/405/406 & route map + request method/Accept & Létezik a route, jó methodot és elfogadható reprezentációt kér a kliens? \\
\addlinespace
409 & optimistic-lock/business state & Közben megváltozott a record, vagy a \code{Changed} query nulla sort módosított? \\
\addlinespace
413/431 & request hard limitek & Body/upload vagy header lépte túl a konfigurált limitet? \\
\addlinespace
421/426 & Host/proxy/TLS policy & Helyes a public host, trusted proxy és az effektív HTTPS? \\
\addlinespace
422 & named form rerender & Melyik mező typed decode/validation hibája áll fenn? \\
\addlinespace
429 & rate-limit metrics/audit & Melyik limiter és melyik identity/IP bucket fogyott el? \\
\addlinespace
500 & request ID alapján server log & Valódi belső invariáns/internal hiba történt? \\
\addlinespace
503 & readiness + server log & DB/Redis/auth/storage vagy request resource limit tette átmenetileg kiszolgálhatatlanná? \\
\addlinespace
DB hiba & DB readiness + query környezet & Kapcsolat, timeout, schema drift, lock vagy constraint? \\
\addlinespace
Fájl/media hiba & AppFs path/jogosultság/limit & Storage root írható, confinement és byte/pixel limit rendben? \\
\addlinespace
Külső API hiba & egress policy + adapter log & DNS, CIDR, port, TLS, timeout vagy response limit zárt? \\
\addlinespace
Log eltűnik rotation után & log fallback metric, service log & SIGHUP eljutott-e, új fájl jogosultsága megfelelő-e? \\
\bottomrule
\end{longtable}

\section{Mit vigyél magaddal productionbe?}

Öt szabály foglalja össze az operátori modellt:

\begin{itemize}
  \item HTTP felület csak typed adat, explicit policy és bounded külső input után legyen;
  \item a stabil policy trusted konfigurációban, a secret pedig secret fájlban maradjon;
  \item schema-módosítás, release-aktiválás és rollback legyen explicit operátori döntés;
  \item proxy, Redis, DB, AppFs és outbound integráció a deklarált authority boundary szerint működjön;
  \item release előtt verifikálj, utána evidence-dzsel bizonyíts healtht, observabilityt, backupot és restore-t.
\end{itemize}

\section{Velran webfejlesztői kompetencia-checklist}

Ez a lista más kérdésre válaszol, mint a production gate: \emph{képes vagy-e egy kisebb Velran webalkalmazást önállóan felépíteni és üzemeltetni anélkül, hogy vakon receptet másolnál?} Egy kipipált pont azt jelentse, hogy a feladatot meg tudod oldani, a releváns boundaryt el tudod magyarázni, és a legvalószínűbb hibát diagnosztizálni tudod.

\subsection*{Alkalmazásfejlesztési alapkompetenciák}
\begin{itemize}[label={\texttt{[ ]}},leftmargin=2em]
  \item El tudok olvasni egy meglévő Velran projektet, és felismerem a page, action, query, route, modul, migration és konfigurációs boundarykat.
  \item Tudok page vagy action handlert hozzáadni, és explicit access policyval route-hoz kötni.
  \item El tudom dönteni, hogy query, form, JSON vagy multipart input illik a feladathoz, és a külső kontroll alatt álló stringeket és collectionöket bounded formában tartom.
  \item Rust-szerű structtal, enummal, \code{Option<T>}, \code{Vec<T>} és \code{Result<T,E>} típussal tudok adatot modellezni legacy Velran írásmód nélkül.
  \item Tudok adatbázis-migrationt hozzáadni, typed queryt írni, és az adatbázis-constraintet összhangban tartani az alkalmazásszintű validációval.
  \item Tudok create, read, update és delete folyamatot írni, és tranzakciót használok, ha több írás egyetlen logikai műveletet alkot.
  \item Felismerem, mikor kell optimistic locking vagy más concurrency szabály a frissebb állapot csendes felülírása helyett.
  \item Egy növekvő alkalmazást modulokra tudok bontani anélkül, hogy gyengíteném a route-, adat- vagy capability-boundarykat.
  \item Tudok typed JSON endpointot készíteni, és külön kezelem a böngészős form-flowt és az API request/response folyamatot.
  \item Tudok adatbázisból érkező Markdownot szerveroldalon renderelni, a Markdownot source of truthként megtartani, és public response cache-t csak valóban nem személyre szabott, megfelelő freshness contracttal rendelkező route-on használok.
  \item A canonical repository-példákat és a Secure Notes checkpointokat használom ellenőrzésre ahelyett, hogy emlékezetből támogatás nélküli szintaxist találnék ki.
\end{itemize}

\subsection*{Security boundaryk}
\begin{itemize}[label={\texttt{[ ]}},leftmargin=2em]
  \item Értem, hogy az authentication azt bizonyítja, ki az actor, az authorization pedig azt dönti el, hogy az actor végrehajthatja-e ezt a műveletet ezen az objektumon.
  \item Tudatosan tudok választani a \code{public}, authenticated, MFA, role, permission, webhook és critical-operation route policyk között.
  \item Védett objektum olvasása vagy módosítása előtt ownership vagy más object-level authorization szabályt érvényesítek.
  \item Principal/tenant szűrést már az adatbázis-querybe teszek, ha egy collection soha nem tartalmazhat más actorhoz tartozó rekordot.
  \item El tudom magyarázni, miért security boundary a korlátlan külső input, upload/media feldolgozás és outbound hálózat, nem egyszerű kényelmi feature.
  \item A CSRF-, CORS-, trusted-proxy-, Host-, cookie- és credential-kezelést külön fogalomként értem, nem egyetlen általános ``web security'' kapcsolóként.
  \item A gyakori compiler/security diagnosticokat értelmezni tudom, és a megsértett invariánst javítom a check elnyomása helyett.
  \item Egy változtatás reviewjakor meg tudom mondani, mely capability-, filesystem-, database-, network- vagy identity-boundary bővül.
\end{itemize}

\subsection*{Build, verification és production}
\begin{itemize}[label={\texttt{[ ]}},leftmargin=2em]
  \item Önállóan végig tudom vinni a normál fejlesztési ciklust: módosítás, \code{velran-cli check}, célzott tesztek, build és a változáshoz tartozó repository verification gate-ek.
  \item A hibakereséshez szükséges mélységben értem a natív útvonalat a verified source-tól az executable IR-en és generált Ruston át a \code{rustc}, immutable generation és atomic activation lépésekig.
  \item Meg tudom különböztetni az alkalmazás reloadolható állapotát azoktól a konfigurációs vagy infrastruktúra-változásoktól, amelyek process restartot igényelnek.
  \item Migrationt elő tudok készíteni és ellenőrizni anélkül, hogy a schema-módosítást a szerver indulásához kötném.
  \item Request ID alapján tudok server/access/audit logot olvasni, és readiness/metrics segítségével különválasztom az alkalmazáshibát a dependency-hibától.
  \item Explicit release artifacttal, checksummal, konfigurációs review-val, smoke teszttel és rollback döntéssel tudok deployolni.
  \item Meg tudom nevezni, ki felel a TLS renewalért, OS patchingért, időszinkronért, lemez/inode kapacitásért, DB/Redis karbantartásért és firewall/DNS változásokért, és tudom, mely Velran checkeket kell ezek után újra lefuttatni.
  \item El tudom mondani, mi tartozik a backupba, izoláltan vissza tudok állítani konzisztens DB/AppFs/identity állapotot, és a restore drillt evidence-nek tekintem, nem puszta dokumentációnak.
  \item Le tudom futtatni a \code{./verify.sh}-t, megértem, melyik gate bukott el, meg tudom különböztetni a tanulói példákat a rejection fixture-öktől/manifesztektől, és kerülöm az olyan ``javítást'', amely csak gyengíti az ellenőrzést.
\end{itemize}

\subsection*{Hogyan használd a checklistet?}
Nem minden parancs bemagolása a cél, hanem az önálló döntés. Ha egy pont bizonytalan, térj vissza a megfelelő Secure Notes checkpointhoz, változtass meg tudatosan egy dolgot, futtasd a releváns verifiert, és mondd el, miért biztonságos az eredmény. Ha a fenti pontok mindegyikét egy kisebb alkalmazáson bizonyítani tudod, elérted azt a gyakorlati alapszintet, amelyet egy Velran webfejlesztőtől elvárunk; a későbbi referenciaanyag ezután mélyít, nem pedig helyettesíti ezt az alapot.

\section{Gyakorlat: evidence-gyűjtésként futtasd a checklistet}
\paragraph{Gyakorlási mód: üzemeltetési szcenárió.} Használd az előszóban meghatározott üzemeltetési módszert.

Ne emlékezetből pipáld ki a pontokat. Minden release-safety szempontból releváns elemhez kapcsolj parancseredményt, config review-t, teszteredményt, restore próbát vagy más konkrét evidence-et.

\section{Tipikus hiba: a checklistet engineering judgment helyettesítőjeként kezeljük}
A checklist ismert omission-osztályokat fog meg; ismeretlen architektúráról nem bizonyítja, hogy biztonságos. Állítsd meg a release-t, ha egy pont kétértelmű, és oldd fel a tervezési kérdést ahelyett, hogy optimistán értelmeznéd a checkboxot.

\section{Végső Secure Notes checkpoint}
E függelékkel végezz release-readiness review-t a végső checkpointon:
\begin{quote}\path{examples/secure-notes/checkpoints/23-release-readiness/main.vrn}\end{quote} Kövesd végig az evidence-láncot a forrástól és lockfile-tól a baseline migráción, compiler verificationön, konfiguráción és natív aktiváláson át a backup/restore eljárásig. Minden nem bizonyítható pont legyen a következő engineering feladat a production aktiválás előtt.
